\chapter{Cybersecurity Act (CSA)}
\label{cap:csa}

\begin{nota}[Fonti]
Lezione del 09/11/2026 (Capitolo~\ref{lez:2026-11-09}). Le slide sul CSA non sono ancora disponibili:
questa è una prima stesura dal parlato; articoli e date da completare sono segnalati.
\end{nota}

\section{Natura e finalità}
Il Cybersecurity Act è un \textbf{regolamento} dell'Unione europea che introduce la \textbf{soft law
europea} della cybersicurezza. Ha due pilastri:
\begin{enumerate}
  \item il \textbf{rafforzamento di ENISA}, con un mandato e compiti propri;
  \item un \textbf{quadro europeo di certificazione} della cybersicurezza per prodotti, servizi e processi ICT.
\end{enumerate}
Nasce per superare quattro problemi:
\begin{itemize}
  \item l'\textbf{interdipendenza}: un fornitore può colpire molti clienti, in più paesi;
  \item l'\textbf{asimmetria informativa} tra venditore e acquirente;
  \item la \textbf{frammentazione} degli schemi nazionali;
  \item la mancanza di \textbf{fiducia verificabile}: servono requisiti ed evidenze.
\end{itemize}
Obiettivo: la \textbf{cyber assurance}, cioè evidenze che certi requisiti sono soddisfatti (da non
confondere con la cyber insurance). Le certificazioni del CSA sono \textbf{volontarie}: oggi in Europa
non esistono certificazioni di cybersicurezza obbligatorie.

\dacompletare{Riferimento normativo (Reg.~UE 2019/881), struttura del regolamento, articoli principali e
mandato permanente di ENISA, dalle slide.}

\section{Hard law, soft law e posizione del CSA}
\begin{center}
\small
\begin{tabularx}{\linewidth}{>{\raggedright\arraybackslash}p{3cm} X X}
\toprule
 & \textbf{Hard law} & \textbf{Soft law}\\
\midrule
Esempi & limiti di velocità; GDPR, CRA, NIS2, DORA, AI Act & policy aziendali (es.\ MFA), linee guida ENISA,
  standard ISO, schemi del CSA\\
Conseguenze & sanzioni predeterminate dalla legge, discrezionalità delimitata & non previste dalla legge,
  ampio margine (es.\ provvedimenti disciplinari per violazione di policy)\\
Tempi & 1--3 anni (o più) per una legge & adattabile in fretta\\
\bottomrule
\end{tabularx}
\end{center}
Il CSA è soft law che aiuta a rispettare la hard law: in particolare, uno schema del CSA può fornire i
requisiti concreti con cui dimostrare la conformità al \textbf{CRA} (Capitolo~\ref{cap:cra}). Offre
\textbf{la possibilità, non la certezza}, di dimostrarla.

\section{ENISA e governance}
\begin{itemize}
  \item \textbf{ENISA}, Agenzia dell'UE per la cybersicurezza: linee guida, formazione ed esercitazioni;
    cooperazione operativa tra CSIRT; \textbf{preparazione degli schemi candidati}; analisi delle minacce e
    buone pratiche.
  \item \textbf{ECCG} (\emph{European Cybersecurity Certification Group}): rappresentanti delle autorità
    nazionali; coopera alla preparazione degli schemi ed esprime un parere sulle bozze.
  \item \textbf{Commissione europea}: individua il bisogno e adotta lo schema.
  \item \textbf{Autorità nazionali di certificazione}: vigilano sulla certificazione; in Italia l'\textbf{ACN}.
  \item \textbf{Organismi di valutazione e laboratori}: eseguono i test e rilasciano i certificati.
\end{itemize}

\section{Processo di adozione di uno schema}
\begin{center}
\begin{tikzpicture}[node distance=4mm, every node/.style={draw, rounded corners, align=center,
    font=\small, minimum height=1.1cm, text width=2.9cm}]
  \node (r) {\textbf{Richiesta}\\Commissione $\rightarrow$ ENISA};
  \node (d) [right=of r] {\textbf{Bozza}\\ENISA + stakeholder};
  \node (o) [right=of d] {\textbf{Parere}\\ECCG};
  \node (a) [right=of o] {\textbf{Adozione}\\Commissione};
  \draw[->] (r) -- (d); \draw[->] (d) -- (o); \draw[->] (o) -- (a);
\end{tikzpicture}
\end{center}
\dacompletare{Atto di adozione (atto di esecuzione), programma di lavoro dell'Unione (URWP), durata e
revisione degli schemi.}

\section{Framework, schema e certificato}
\begin{itemize}
  \item \textbf{Framework}: il CSA, che fissa ruoli e regole comuni. Non contiene schemi.
  \item \textbf{Schema}: il \emph{rule book} della valutazione per una categoria: ambito, requisiti di sicurezza,
    test, livelli di affidabilità.
  \item \textbf{Certificato}: attesta che uno specifico prodotto o servizio rispetta lo schema.
\end{itemize}
Oggetti possibili di uno schema: \textbf{prodotti ICT} (firewall, chip, software), \textbf{servizi ICT} (es.\ cloud
storage), \textbf{processi ICT} (progettazione e manutenzione; ancora nessuno schema), \textbf{servizi di
sicurezza gestiti} (incident response, audit, penetration test). ISO/IEC~27001 certifica invece
l'\textbf{organizzazione}.

\section{Livelli di affidabilità e valutazione della conformità}
\begin{center}
\small
\begin{tabular}{lll}
\toprule
\textbf{Livello} & \textbf{Rischio} & \textbf{Valutazione}\\
\midrule
Base & rischi noti di base & anche autovalutazione del fabbricante\\
Sostanziale & attaccanti con risorse limitate & terza parte\\
Elevato & attaccanti con risorse significative & terza parte\\
\bottomrule
\end{tabular}
\end{center}
In pratica la certificazione è quasi sempre di terza parte. Un certificato:
\begin{itemize}
  \item dà evidenze su un oggetto \textbf{identificato con precisione}, su requisiti e test specifici, a un
    livello definito;
  \item \textbf{non} dimostra l'assenza di vulnerabilità;
  \item non esonera in automatico da responsabilità: dopo un incidente conta se i requisiti erano davvero
    rispettati (\emph{accountability});
  \item richiede \textbf{monitoraggio e manutenzione}: prodotto e minacce cambiano; una patch non invalida
    automaticamente il certificato, ma il suo impatto va valutato; una nuova versione va ricertificata.
\end{itemize}
\dacompletare{Articoli su livelli di affidabilità, autovalutazione della conformità, durata e revoca dei
certificati, sanzioni nazionali.}

\section{EUCC}
Primo schema adottato: \textbf{EUCC} (\emph{EU Common Criteria-based scheme}), basato sui Common Criteria;
riguarda \textbf{prodotti e componenti ICT} (niente servizi o processi), è applicabile dal
\textbf{27 febbraio 2025} e prevede i livelli \textbf{sostanziale ed elevato}.
Concetti: \textbf{target of evaluation} (il prodotto esatto), \textbf{security target} (il documento che lo
descrive con i requisiti dichiarati), \textbf{protection profile} (requisiti riutilizzabili per un tipo di
prodotto), metodi e analisi delle vulnerabilità, ipotesi e limiti operativi.
\dacompletare{Riferimento al regolamento di esecuzione EUCC; altri schemi in preparazione (EUCS per il
cloud, EU5G).}

\section{Rapporto con le altre norme e strumenti}
\begin{center}
\small
\begin{tabularx}{\linewidth}{>{\raggedright\arraybackslash}p{3.4cm} X}
\toprule
\textbf{Norma / strumento} & \textbf{Domanda a cui risponde}\\
\midrule
CSA & come si valutano le dichiarazioni di sicurezza: quadro comune di certificazione\\
CRA & obblighi di sicurezza e sulle vulnerabilità dei prodotti con elementi digitali\\
NIS2 & sicurezza di reti e sistemi informativi dei soggetti essenziali e importanti\\
GDPR & protezione dei dati personali (la norma più generale)\\
DORA & resilienza operativa digitale del settore finanziario\\
ISO/IEC~27001 & sistema di gestione della sicurezza dell'organizzazione\\
Marcatura CE & conformità del prodotto alle norme UE che la richiedono (possibile sovrapposizione con l'EUCC)\\
Penetration test & attacco simulato autorizzato; soft law con possibili conseguenze legali\\
\bottomrule
\end{tabularx}
\end{center}

\section{Security by design e supply chain}
Protezione e manutenzione si pianificano in progettazione: ridurre la superficie d'attacco, conoscere i
componenti (\textbf{SBOM}), gestire le segnalazioni di vulnerabilità, distribuire gli aggiornamenti. Per ogni
fornitore: il componente può essere attaccato? Possiamo operare senza? Chi controlla servizi, decisioni e
dati (governance, non solo \emph{data residency})?

\section{CSA~2}
Proposta di revisione: più poteri a ENISA, certificazione semplificata, più attenzione alla supply chain.
\dacompletare{Stato dell'iter legislativo e contenuti della proposta.}

\begin{esame}
\begin{itemize}
  \item I due pilastri: \textbf{ENISA} e \textbf{schemi di certificazione}.
  \item CSA = soft law sui prodotti; ISO/IEC~27001 = organizzazione; CRA = hard law sui prodotti.
  \item Framework $\neq$ schema $\neq$ certificato; un certificato non prova l'assenza di vulnerabilità.
  \item Livelli base/sostanziale/elevato; autovalutazione solo per il livello base.
  \item EUCC: primo schema, prodotti ICT, Common Criteria, livelli sostanziale ed elevato.
\end{itemize}
\end{esame}
